% Discussion for item3_1 -- prose goes here.
% scripts/make_latex.py will NEVER overwrite this file.

\textbf{Address of Pointer} (\texttt{\&a}, \texttt{\&b}) reflects the local
pointer variables on the stack, separated by 8 bytes, the pointer size on this
64-bit architecture.

\textbf{Effective Address} prints the values held in \texttt{a} and \texttt{b}
before any allocation. Reading them is undefined behaviour: they hold whatever
the stack happened to contain. In this run one evaluated to a nonsensical
address and the other to \texttt{(nil)}, which is precisely the point ---
\texttt{(nil)} is not guaranteed, and an uninitialised pointer is not a null
pointer.

\textbf{After malloc} shows \texttt{a} pointing to dynamic memory on the heap,
far from the stack. \texttt{a} and \texttt{\&a[0]} reference the start of the
buffer, while \texttt{\&a[9]} is offset by 36 bytes ($9 \times 4$ bytes per
\texttt{int}).

\textbf{Array c} is allocated on the stack alongside the other automatic
variables, requiring no dynamic allocation and permitting no explicit
deallocation; its lifetime ends with the function.

\textbf{After realloc} preserves \texttt{a}'s base address, because no
allocation was made immediately after it and the block could be extended in
place, up to \texttt{\&a[999]} at 3996 bytes past the base.
